% =============================================================================
% Seção 8 — Conclusões e trabalhos futuros
% =============================================================================
\section{Conclusões e trabalhos futuros}
\label{sec:conclusoes}

Este artigo apresentou uma abordagem computacional para o pré-dimensionamento geométrico de fundações superficiais do tipo sapata isolada, formulada como otimização mono-objetivo penalizada e resolvida por uma arquitetura híbrida \textit{Efficient Global Optimization} (EGO) com modelo substituto por Regressão por Processos Gaussianos (GPR) e Algoritmo Genético (AG) como otimizador interno da função de aquisição \textit{Expected Improvement} (EI). O recorte corresponde à primeira frente metodológica do plano de trabalho — dimensões de sapatas com posições de pilares fornecidas pelo projeto estrutural e verificações parciais de tensão, geometria e punção —, agora sustentado por um protocolo experimental fechado: 30 repetições semeadas por célula, orçamento de avaliações equalizado, dois cenários de comparação, teste estatístico pareado e artefatos integralmente persistidos e regeneráveis.

As conclusões consolidadas são as seguintes:
\begin{enumerate}
    \item \textbf{Eficiência amostral demonstrada nas instâncias avaliadas.} Sob orçamento igual de 150 avaliações reais, as arquiteturas assistidas por modelo substituto ocupam o topo em $\Theta$ nos três casos. A redução mediana da melhor arquitetura assistida em relação à busca aleatória foi de $24{,}7\,\%$, $45{,}7\,\%$ e $71{,}4\,\%$ na sequência de casos testada; para o EGO isolado, foi de $23{,}4\,\%$, $38{,}8\,\%$ e $63{,}9\,\%$. Como há uma instância por dimensionalidade e a restrição de sobreposição permanece inativa, esse padrão não deve ser interpretado como isolamento causal do efeito da dimensão.
    \item \textbf{A quase separabilidade foi quantificada, não apenas declarada.} A linha de base por Differential Evolution por sapata, seguida de reavaliação global estrita, encontrou volumes factíveis de \SI{3,109}{\meter\cubed}, \SI{4,751}{\meter\cubed} e \SI{2,122}{\meter\cubed}. Esses valores empatam ou melhoram o melhor volume factível obtido no protocolo global em $<0{,}01\,\%$, $0{,}77\,\%$ e $2{,}08\,\%$, respectivamente. Portanto, os casos atuais são fortes para validar implementação, métricas, penalização e CBO, mas ainda fracos para demonstrar comportamento em problemas acoplados.
    \item \textbf{O tratamento explícito de restrições melhora $\Theta$, mas exige seleção factível.} A \textit{Constrained Bayesian Optimization} (CBO), que modela volume e restrições por modelos substitutos independentes com aquisição por probabilidade de factibilidade \citepabnt{gardner2014bayesian}, foi confrontada com o EGO penalizado sob sementes, orçamento e demais parâmetros idênticos. O CBO reduziu a média de $\Theta$ em $1{,}5\,\%$, $9{,}3\,\%$ e $21{,}7\,\%$ frente ao EGO nos três casos, e melhorou o melhor volume estritamente factível em $0{,}8\,\%$, $3{,}5\,\%$ e $15{,}3\,\%$. Em contrapartida, sua factibilidade estrita caiu para $63\,\%$ e $37\,\%$ nos casos 1 e 2, permanecendo em $83\,\%$ no caso 3. A variante escalável \citepabnt{eriksson2021scbo} permanece como hipótese prioritária para a etapa realmente acoplada de empacotamento, acompanhada por regras de seleção final factível.
    \item \textbf{Domínio de adequação delimitado com precisão.} Com a função pseudo-objetivo vetorizada ($\approx \SI{0,1}{\milli\second}$ por avaliação), buscas diretas com orçamento estendido ($3\,000$ avaliações) superam as referências assistidas de 150 avaliações em qualidade consumindo cerca de $1\,\%$ do tempo de parede. As arquiteturas assistidas são, portanto, adequadas quando avaliações reais são o recurso escasso — o regime previsto para as extensões da pesquisa com empacotamento, interação geotécnica e recalques —, e não uma necessidade computacional da formulação algébrica atual. Essa delimitação empírica substitui a justificativa genérica de custo usada em etapas preliminares.
    \item \textbf{Penalidade governa o modelo substituto; o \textit{kernel} foi secundário no rastreio.} O estudo controlado kernels~$\times$~penalidade mostrou que o $R^2$ global é insensível ao fator de penalidade (a variância artificial da penalização é reproduzida pelo GPR), mas o erro absoluto na região factível — onde a aquisição decide — cresce cinco ordens de grandeza sob $\alpha = 10^{6}$, fundamentando a penalidade moderada adotada. Com apenas três réplicas por configuração de \textit{kernel}, a escolha fina da covariância deve ser lida como triagem exploratória; a configuração de produção (Matérn $\nu = 2{,}5$) permaneceu na banda superior ($R^2 = 0{,}931 \pm 0{,}011$).
    \item \textbf{Comportamento previsto da penalização exterior, agora quantificado.} Com $\alpha = 10$ e $p = 1$, soluções finais de todos os métodos carregam violações residuais ($g_k \sim 10^{-4}$--$10^{-1}$; factibilidade estrita de $30$ a $83\,\%$). Essas violações não são aceitáveis como soluções de projeto; elas indicam que a seleção final deve priorizar factibilidade estrita ou regras de dominância de factibilidade. A tensão admissível do solo é a restrição governante nos três casos; a punção — verificada nos dois contornos críticos da NBR~6118, $C$ e $C'$ — permanece folgada na faixa explorada; e a restrição de sobreposição é inativa por construção nos casos congelados.
    \item \textbf{Implementação consolidada e reprodutível.} A plataforma \fundaai{} integra leitura de dados, validações de fronteira da formulação (carga vertical estritamente positiva, plausibilidade de unidades do $f_{ck}$, altura útil positiva em todo o espaço de busca), execução dos métodos, bancada de comparação com orçamento controlado, persistência autodescritiva de execuções e exportação (Excel, DXF, JSON), viabilizando a regeneração integral dos resultados deste artigo por \textit{scripts} determinísticos.
\end{enumerate}

\subsection{Trabalhos futuros}
\label{sec:trabalhos_futuros}

\paragraph{Frente 1 -- Posicionamento conjunto e empacotamento.}
A próxima etapa central consiste em incorporar o problema de empacotamento ao processo de otimização, tratando as posições (ou excentricidades) das sapatas como variáveis de projeto sob restrições rígidas de não sobreposição, fronteira do lote e momentos adicionais induzidos. Nesse regime, a otimização deixa de ser quase separável por sapata e passa a ter acoplamento geométrico real. Essa frente exigirá revisão bibliográfica específica de empacotamento bidimensional e \textit{layout optimization}.

\paragraph{Frente 2 -- Aprofundamento do tratamento de restrições.}
A \textit{Constrained Bayesian Optimization} já implementada nesta pesquisa (Seção~\ref{sec:met_cbo}) na formulação de \textit{Expected Constrained Improvement} \citepabnt{gardner2014bayesian} abre três desdobramentos naturais: (i) variantes escaláveis com regiões de confiança \citepabnt{eriksson2021scbo} para o regime de maior dimensionalidade da frente de empacotamento; (ii) aquisições que ponderam explicitamente o risco de infactibilidade quando a base amostral é curta, mitigando o modo de falha observado no caso de fronteira ativa; e (iii) comparação com penalização adaptativa e regras de dominância de factibilidade no laço evolutivo. Para uma etapa posterior, modelos substitutos alternativos para CBO, como redes pré-treinadas do tipo PFN avaliadas por \citeonline{yu2025fast}, podem ser investigados apenas se o custo de ajuste de múltiplos GPs se tornar gargalo efetivo. A linha de base estatística deste artigo permite julgar cada variante com o mesmo protocolo.

\paragraph{Frente 3 -- Validação ampliada e incerteza.}
A consolidação adicional contempla: (i) inclusão de CMA-ES, Differential Evolution e otimizadores determinísticos como métodos pareados no mesmo protocolo de orçamento; (ii) casos com maior número de elementos e o dimensionamento explícito da armadura de flexão, com verificação de cisalhamento unidirecional, ancoragem e custo total em lugar de apenas volume de concreto, aproximando o escopo de estudos de projeto econômico completo \citepabnt{nigdeli2018metaheuristic}; (iii) reconstrução formal das combinações de ações para separar verificações de serviço e de estado limite último; (iv) validação contra exemplos resolvidos da bibliografia clássica de fundações; e (v) extensão para incerteza geotécnica e de carregamento, na linha de projeto robusto de \citeonline{juang2013reliability}, incorporando recalques e modelos de capacidade do solo calibrados quando houver base geotécnica suficiente \citepabnttwo{khajehzadeh2022effective}{fattahi2025settlement}.

\subsection{Considerações finais}

Os resultados aqui reportados sustentam uma conclusão metodológica delimitada: para o pré-dimensionamento geométrico de sapatas isoladas com posições fixas e orçamento restrito de avaliações, as arquiteturas assistidas por modelo substituto foram competitivas ou superiores às buscas diretas testadas em $\Theta$; enquanto a avaliação permanece de baixo custo, porém, buscas diretas com orçamento maior e a própria decomposição por sapata aproximam melhor o ótimo factível. As reduções de volume penalizado frente às buscas não assistidas, a caracterização quantitativa do efeito da penalidade sobre o modelo substituto, a auditoria de quase separabilidade e a delimitação explícita de factibilidade e escopo convertem a plataforma \fundaai{} em base experimental para a frente de empacotamento, na qual o posicionamento conjunto das fundações ativará restrições acopladas e exigirá nova validação.
